%% kapa: exam=1 round=33 year=2022 period=2 form=A subject=accounting subject_ko=회계학 seq=7 number=47
%% source: Q-Net 감정평가사 기출문제 (cst00309) artlSeq=5209997 file=raw/1008/5209997/2022년도 제33회 감정평가사 1차시험 2교시(감정평가 관계법규, 회계학).pdf
%% conversion: kordoc 4.24.51 -> tools/gen1.py (auto); review=proofread-sonnet; flags=none
\begin{question}{47}
재무제표 요소에 관한 설명으로 옳지 않은 것은?
\begin{choices}
  \item 자산은 과거사건의 결과로 기업이 통제하는 현재의 경제적자원이다.
  \item 부채는 과거사건의 결과로 기업이 경제적자원을 이전해야 하는 현재의무이다.
  \item 수익은 자본청구권 보유자로부터의 출자를 포함하며, 자본청구권 보유자에 대한 분배는 비용으로 인식한다.
  \item 기업이 발행한 후 재매입하여 보유하고 있는 채무상품이나 지분상품은 기업의 경제적자원이 아니다.
  \item 자본청구권은 기업의 자산에서 모든 부채를 차감한 후의 잔여지분에 대한 청구권이다.
\end{choices}
\end{question}
\answer{3}
